\documentclass{article}
\usepackage[utf8]{inputenc}
\usepackage{geometry}
\usepackage{amsmath}
\usepackage{multicol}

\geometry{a4paper, margin=1in}

\title{Plan de Análisis: Modelado Multi-Agente (MARL) para Simulación de Conductores por Clústeres}
\author{Equipo de Desarrollo IA - SUMO}
\date{\today}

\begin{document}

\maketitle

\section{Introducción}
Este documento detalla el plan de análisis para el desarrollo de un sistema de Aprendizaje por Refuerzo Multi-Agente (MARL). El objetivo principal es entrenar múltiples agentes inteligentes, donde cada agente representa y aprende el comportamiento de conducción de un clúster específico de conductores, identificado previamente en la fase de \textit{Clustering}. El entrenamiento se realizará utilizando datos históricos de flujo vehicular, permitiendo simular escenarios de tráfico realistas y heterogéneos en SUMO.

\section{Definición del Problema}

El problema consiste en resolver una tarea de \textit{Imitación o Calibración Dinámica} mediante RL. Dado un conjunto de datos de flujo $D_{real}$ que contiene métricas agregadas (velocidad, flujo, densidad) desglosadas o inferidas por clúster, buscamos una política $\pi_c$ para cada clúster $c \in C$ tal que el tráfico simulado $D_{sim}$ generado por los agentes se asemeje estocásticamente a $D_{real}$.

\subsection{Entradas del Sistema}
\begin{itemize}
    \item \textbf{Definición de Clústeres}: Perfiles de comportamiento (e.g., prudente, agresivo, errático) obtenidos de `src/clustering.py`.
    \item \textbf{Datos de Flujo (Snapshot Features)}: Series temporales de velocidad, flujo y densidad por pórtico y carril, procesados en `src/snapshot\_pipeline.py`.
    \item \textbf{Topología del Grafo}: Red vial definida en SUMO y mapeada en el sistema GNN.
\end{itemize}

\section{Modelado MARL}

\subsection{Agentes}
Se definirá un conjunto de agentes $A = \{a_1, ..., a_N\}$, donde $N$ es el número de clústeres.
\begin{itemize}
    \item Un agente no representa un vehículo individual, sino un \textbf{controlador de tipo de conductor}.
    \item El agente $a_c$ controla los parámetros de comportamiento (o acciones directas) de todos los vehículos asignados al clúster $c$ en la simulación.
\end{itemize}

\subsection{Espacio de Estados ($S$)}
El estado $s_t$ debe capturar la dinámica actual del tráfico para que el agente tome decisiones informadas. Basado en `src/snapshot\_pipeline.py`, el estado incluirá:
\begin{itemize}
    \item \textbf{Estados Locales}: Velocidad media, densidad y flujo actual en los ejes de interés.
    \item \textbf{Contexto Temporal}: Features de ventana (medias móviles, desviaciones estándar) de los últimos $k$ pasos.
    \item \textbf{Contexto Espacial}: Gradientes de velocidad/densidad con tramos aguas arriba y aguas abajo.
    \item \textbf{Estado del Clúster}: Métricas actuales de desempeño del propio clúster en la simulación (e.g., "mis vehículos van muy lento comparado con el real").
\end{itemize}

\subsection{Espacio de Acciones ($A$)}
Las acciones $u_t^c$ para el agente del clúster $c$ pueden ser:
\begin{itemize}
    \item \textbf{Ajuste de Parámetros de Car-Following (Continuo)}: Modificar dinámicamente parámetros del modelo de Krauss/IDM en SUMO (e.g., `maxSpeed`, `accel`, `decel`, `minGap`, `sigma`).
    \item \textbf{Ajuste de Cambio de Carril (Continuo/Discreto)}: Parámetros de `lcStrategic`, `lcCooperative`.
    \item \textit{Nota}: Las acciones se aplican globalmente a los vehículos del clúster $c$ en el instante $t$ o a un subconjunto relevante.
\end{itemize}

\subsection{Función de Recompensa ($R$)}
El objetivo es minimizar la discrepancia entre simulación y realidad. La recompensa $r_t^c$ para el agente $c$ se define como:
\begin{equation}
    r_t^c = - \left( w_v \| v_{sim}^c - v_{real}^c \|^2 + w_f \| f_{sim}^c - f_{real}^c \|^2 + w_{safe} \cdot \text{Penalty}_{crash} \right)
\end{equation}
Donde:
\begin{itemize}
    \item $v_{sim}^c, v_{real}^c$: Velocidades medias del clúster $c$.
    \item $f_{sim}^c, f_{real}^c$: Flujos/Densidades del clúster $c$.
    \item $\text{Penalty}_{crash}$: Penalización por colisiones excesivas (restricción de seguridad).
\end{itemize}

\section{Restricciones y Desafíos}
\begin{enumerate}
    \item \textbf{Asignación de Clústeres en Datos Reales}: Los datos de flujo agregados (pórticos) no siempre distinguen clústeres. Se requerirá un paso de \textit{inferencia} o desagregación basado en las proporciones estimadas (e.g., mezcla de gaussianas).
    \item \textbf{Seguridad vs. Realismo}: Agentes "agresivos" podrían aprender a chocar para igualar velocidades altas. Se deben imponer "Safety Shields" o penalizaciones duras por colisiones.
    \item \textbf{Estabilidad de SUMO}: Cambios bruscos en parámetros de car-following pueden causar inestabilidades físicas en la simulación.
\end{enumerate}

\section{Plan de Análisis y Ejecución}
\begin{enumerate}
    \item \textbf{Análisis de Distribución de Clústeres}: Determinar las proporciones de cada clúster en los pórticos usando `clustering.py`.
    \item \textbf{Diseño de Entorno Gym-SUMO}: Crear un entorno compatible con OpenAI Gym/PettingZoo que interactúe con SUMO (usando TraCI/Libsumo).
    \item \textbf{Implementación de Algoritmo}: Seleccionar PPO (Proximal Policy Optimization) multi-agente con \textit{Centralized Training, Decentralized Execution} (CTDE) si es posible, o PPO independiente (IPPO) como baseline.
    \item \textbf{Validación}: Comparar histogramas de velocidad y headways simulados vs reales.
\end{enumerate}

\end{document}
